\begin{abstract}
Large-scale unsupervised language models (LMs) are effective at acquiring general world knowledge and some reasoning capabilities, yet exerting fine-grained control over their outputs remains challenging due to the fully unsupervised nature of their training objectives. To impart this controllability, prevailing approaches leverage human-provided labels that express preferences between model outputs, subsequently fine-tuning the original LM to reflect these choices—most commonly using reinforcement learning from human feedback (RLHF). Nevertheless, RLHF procedures are typically complex and can be unstable; they require first learning a reward model that captures the human preference data, followed by reinforcement learning to optimize the unsupervised LM for this reward while constraining deviations from the pre-trained distribution. In this work, we propose a novel formulation of the reward model within RLHF that permits analytical derivation of the associated optimal policy. This allows us to resolve the RLHF challenge via a straightforward classification loss alone. Our approach, termed \textit{Direct Preference Optimization} (DPO), is robust, effective, and highly efficient, removing the necessity for sampling from the LM during adaptation or performing extensive hyperparameter searches. Experimental results indicate that DPO achieves preference alignment on par with or better than conventional techniques. Specifically, DPO fine-tuning offers superior sentiment control compared to RLHF with PPO, and delivers equal or improved performance in response quality for summarization and single-turn dialogue tasks, all while offering significant implementation and training simplicity.
\end{abstract}

\section{Introduction}
Large-scale unsupervised language models (LMs), trained on extensive datasets, have demonstrated unexpected proficiencies~\citep{chowdhery2022palm, brown2020language, touvron2023llama,bubeck2023sparks}. Nonetheless, the training data for these models originates from a diverse range of human contributors, each possessing different objectives, skill levels, and priorities. Not all of these human traits are suitable targets for emulation; for instance, while an AI coding assistant ought to \textit{recognize} frequent programming errors to facilitate their correction, during code generation we prefer the model to emulate the exceptional, if uncommon, high-quality coding instances in its corpus. Likewise, although it is beneficial for the language model to be \textit{informed} of widespread misconceptions held by, say, half of the population, we explicitly do not want the model to endorse such misconceptions in half of its responses to related queries. In essence, it is essential to distinguish and promote the model’s \emph{intended outputs and conduct} from its broader \textit{reservoir of knowledge and skills} in order to develop AI systems that are reliable, effective, and subject to human direction~\citep{ouyang2022training}. Traditionally, guiding LMs toward preferred behaviors has relied on reinforcement learning (RL) to align their outputs with human judgments; yet, as we will demonstrate, the prevailing RL-based formulation used in these approaches can be precisely solved using a simple binary cross-entropy loss, thereby streamlining the preference alignment process significantly.

Broadly speaking, contemporary strategies inject preferred conduct into LMs using meticulously assembled collections of human preference data, characterizing responses deemed safe and advantageous by human evaluators. This step of preference alignment is performed subsequent to a period of unsupervised pre-training over expansive textual data. Although preference learning could be implemented through direct supervised fine-tuning on exemplary human-provided responses, the most effective strategies—namely reinforcement learning from human or artificial feedback (RLHF/RLAIF; \citep{christiano2017deep,bai2022constitutional})—employ reward modeling to map human preference data, followed by policy optimization via RL to nudge the model toward outputs that score highly according to these rewards, all while maintaining proximity to the pre-trained distribution. While RLHF has enabled the development of models with outstanding dialogue and programming performance, its adoption necessitates a far more involved pipeline than supervised methods, involving the coordination and training of multiple LMs and frequent sampling from the model within the RL loop, which substantially increases computational demands.

This work presents a method for directly optimizing a language model according to human preferences, foregoing the need for explicit reward modeling or the application of reinforcement learning. We introduce \textit{Direct Preference Optimization (DPO)}, a training procedure that inherently maximizes the same objective targeted by typical RLHF techniques (i.e., reward maximization subject to a KL-divergence penalty), yet is markedly straightforward in terms of both implementation and training. At a high level, DPO adjusts model parameters to favor preferred responses over less favored ones by increasing their relative log probabilities, but incorporates dynamically computed, example-specific importance weights that address issues of model collapse seen with naive probability-ratio-based objectives. DPO, akin to other approaches, utilizes a theoretical preference framework (e.g., the Bradley-Terry model; \cite{bradley1952rankanalysis}) to evaluate how well a reward signal corresponds to observed preference judgements. In contrast to conventional approaches, which employ this preference framework to construct a loss function for reward model training and then derive a policy via reward model optimization, DPO leverages a change-of-variables technique to express the preference loss directly in terms of the policy parameters. As a result, given annotated pairwise preference data over model outputs, DPO enables direct policy optimization via a straightforward binary cross-entropy loss, yielding the optimal policy relative to a latent reward function that is consistent with the provided preferences.

The principal innovation of this study is the Direct Preference Optimization (DPO) method—a reinforcement learning-free approach for training language models on preference data. Empirical results demonstrate that DPO matches or exceeds the performance of current methods, including PPO-based RLHF, across diverse tasks such as sentiment adjustment, text summarization, and conversational response generation, on models scaling up to 6B parameters.

\section{Related Work}

Large self-supervised language models are increasingly capable of performing certain tasks either without any additional training data (zero-shot) \citep{radford2019language} or when provided with a few examples as prompts \citep{gpt3,megatron,chowdhery2022palm}. Nevertheless, their effectiveness on various downstream applications and their ability to align with user objectives are markedly enhanced through fine-tuning with collections of instructions paired with human-generated responses \citep{mishra-etal-2022-cross,sanh2022multitask,chung2022scaling,thoppilan2022lamda}. This process, known as `instruction-tuning', allows LLMs to better extrapolate to unseen instructions beyond those present in the instruction-tuning corpus, leading to generally improved model utility \citep{chung2022scaling}. Although instruction tuning has demonstrated notable progress, collecting human preference data—comparative evaluations of responses—tends to be simpler than obtaining detailed expert demonstrations. Consequently, later research has further refined LLMs by training them on datasets reflecting human preferences, which has led to improved outcomes in areas like translation \citep{kreutzer-etal-2018-reliability}, summarization \citep{stiennon2022learning,ziegler2020finetuning}, narrative generation \citep{ziegler2020finetuning}, and instruction adherence \citep{ouyang2022training,ramamurthy2023is}. The typical approach involves initially learning a reward model—a neural network trained to fit preference data under frameworks such as the Bradley-Terry model \citep{bradley1952rankanalysis}—and subsequently optimizing the language model to maximize this reward signal with reinforcement learning methods, including REINFORCE \citep{williams1992reinforce}, proximal policy optimization (PPO; \cite{schulman2017proximal}), or related algorithms \citep{ramamurthy2023is}. In parallel, related studies have utilized LLMs, refined via instruction-following and human feedback, to generate synthetic preference datasets targeting specific characteristics like safety or harmlessness \citep{bai2022constitutional}, relying on limited direct human supervision by applying rubric-based guidelines to the LLM’s outputs. These efforts represent the intersection of two research traditions: one focused on reinforcement learning-based optimization of language models for diverse objectives~\citep{Ranzato2015SequenceLT,paulus2018a,wu2018learning}, and the other on designing methods for learning from human preferences \citep{christiano2017deep,kupcsik2018learning}. Despite the intuitive advantages of leveraging relative human feedback, the practicalities of applying reinforcement learning to large language models remain daunting. This work introduces a theoretically grounded method for directly optimizing relative preferences, thereby bypassing the necessity for reinforcement learning.

Beyond linguistic applications, the study of learning policies from preference information has been extensively explored within both bandit and reinforcement learning frameworks, with a range of methodologies introduced in the literature. Contextual bandit formulations that rely on preferences or rankings, rather than direct reward feedback, are referred to as contextual dueling bandits (CDB; \cite{yue2012karmed,dudik2015contextual}). When absolute rewards are unavailable, the theoretical treatment of CDBs replaces the classical concept of an optimal policy with a \textit{von Neumann winner}, characterized as a policy that, in expectation, secures a win rate of at least 50\% against \textit{any} competing policy \citep{dudik2015contextual}. A notable distinction of the CDB setting is that preference annotations are provided in an online fashion, whereas in typical human preference learning, algorithms operate on a predetermined, static collection of offline preference-labeled action pairs \citep{yan2022human}. Likewise, \textit{preference-based RL} (PbRL) learns via binary preferences derived from an \textit{unknown} underlying ‘scoring’ function as opposed to explicit rewards \citep{BusaFekete2014,ruiz2023dueling}. A multitude of PbRL algorithms have been proposed, including those capable of leveraging off-policy preference datasets, but they generally involve an explicit estimation of the latent scoring function (that is, learning a reward model) prior to policy optimization \citep{jain2013learning,BusaFekete2014,christiano2017deep,sadigh2017active,kupcsik2018learning}. In contrast, our contribution lies in proposing a unified approach in which the policy is trained in a single stage to directly conform to specified preferences.

\section{Preliminaries}\label{section:prelims}

We briefly summarize the RLHF framework as described in \citeauthor{ziegler2020finetuning} and subsequent works (\citep{stiennon2022learning, bai2022training, ouyang2022training}). This process generally comprises three sequential stages: (1) supervised fine-tuning (SFT), (2) collection of preference data and reward modeling, and (3) reinforcement learning-based policy optimization.

\textbf{SFT}: The initial phase of RLHF involves refining a pre-trained language model through supervised learning on a curated dataset that reflects the target downstream tasks (such as dialogue generation, summarization, etc.), resulting in the supervised fine-tuned model $\pi^\text{SFT}$.

\textbf{Reward Modelling Phase}: In the subsequent stage, prompts $x$ are provided to the SFT model, which generates answer pairs $(y_1, y_2)\sim \pi^\text{SFT}(y \mid x)$. These responses are evaluated by human annotators, who indicate a preference between the two, noted as $y_w\succ y_l \mid x$, where $y_w$ and $y_l$ denote the selected (preferred) and unselected (less preferred) completions from the pair, respectively. The source of these preferences is presumed to be some latent reward function $r^*(y, x)$, which remains unobserved. Multiple strategies are available for modeling preferences; the Bradley-Terry (BT) \cite{bradley1952rankanalysis} model is a standard choice (though more general frameworks, such as Plackett-Luce ranking models \citep{plackett1975analysis, luce2012individual}, can also be incorporated when multiple ranked completions are available). Under the BT model, the human preference probability distribution $p^*$ is defined as:

\begin{equation}\label{eq:bradley-terry}
    p^*(y_1\succ y_2 \mid x)=\frac{\exp\left(r^*(x, y_1)\right)}{\exp\left(r^*(x, y_1)\right) + \exp\left(r^*(x, y_2)\right)}.
\end{equation}

Suppose we are provided with a fixed dataset of pairwise comparisons, denoted as $\mathcal{D}=\bigl\{x^{(i)}, y_w^{(i)}, y_l^{(i)}\bigr\}_{i=1}^N$, drawn according to $p^*$. We can then introduce a parameterized reward function $r_{\phi}(x, y)$ and fit its parameters using maximum likelihood estimation. By treating the learning task as a binary classification, we define the corresponding negative log-likelihood objective as:

\begin{equation}\label{eq:reward_model}
    \mathcal{L}_R(r_{\phi}, \mathcal{D}) = -\mathbb{E}_{(x, y_w, y_l)\sim \mathcal{D}}\bigl[\log \sigma(r_{\phi}(x, y_w)- r_{\phi}(x, y_l))\bigr]
\end{equation}

where $\sigma$ denotes the logistic sigmoid function. Within the setting of language models, $r_{\phi}(x, y)$ is typically initialized by extending the SFT model $\pi^\text{SFT}(y \mid x)$ with an added linear head on top of the last transformer block, yielding a scalar reward estimate~\cite{ziegler2020finetuning}. To stabilize training and reduce reward variance, previous approaches have normalized the rewards so that $\mathbb{E}_{x,y\sim \mathcal{D}}\left[r_\phi(x, y)\right] = 0$ for all $x$.

\textbf{RL Fine-Tuning Phase}: In the reinforcement learning phase, the reward model provides signal to the language model policy. As in earlier studies~\citep{jaques2017sequence, jaques2020human}, the objective to be optimized is:

\begin{equation}\label{eq:RL}
\max_{\pi_{\theta}}  \mathbb{E}_{x\sim \mathcal{D}, y\sim \pi_{\theta}(y \mid x)}\bigl[r_{\phi}(x, y)\bigr] - \beta\mathbb{D}_{\textrm{KL}}\bigl[\pi_{\theta}(y\mid x)\mid \mid \pi_\text{ref}(y\mid x)\bigr],
\end{equation}

where $\beta$ determines how strongly the policy is regularized to remain close to the reference policy $\pi_\text{ref}$, which is typically the initial SFT model, $\pi^\text{SFT}$. 
In most implementations, the policy $\pi_\theta$ starts from $\pi^\text{SFT}$ as well. This regularization is critical for preventing excessive drift from the domain where the reward model is valid, preserving sample diversity, and avoiding degeneration into narrow, repetitive outputs. Owing to the discrete property of language outputs, direct gradient-based optimization of this objective is infeasible, so reinforcement learning methods are usually employed. Conventionally~\citep{ziegler2020finetuning, stiennon2022learning, bai2022training, ouyang2022training}, the reward is defined as ${r(x, y) = r_{\phi}(x, y) -\beta (\log \pi_{\theta}(y\mid x) - \log \pi_\text{ref}(y\mid x))}$, and optimization is conducted with PPO~\cite{schulman2017proximal}. 

\section{Direct Preference Optimization}\label{sec:DPO}

Given the complexities inherent to reinforcement learning for large-scale tasks such as language model fine-tuning, our objective is to establish a streamlined method for optimizing policies directly using preference data. In contrast with standard RLHF pipelines—which entail reward model learning followed by reinforcement learning optimization—we propose exploiting a specific form of reward function parameterization that allows the optimal policy to be recovered analytically, removing the need for a reinforcement learning procedure. As detailed below, our central contribution is the identification of an analytic correspondence from reward functions to optimal policies; this facilitates reexpressing reward-based objectives as direct policy-based objectives through a change of variables. With this strategy, we sidestep the explicit training of an isolated reward model, while continuing to train under established preference modeling frameworks such as the Bradley-Terry model. Conceptually, the policy

\textbf{Deriving the DPO objective.} We begin by considering the general reinforcement learning objective expressed in Eq.~\ref{eq:RL}, incorporating an arbitrary reward function $r$. Drawing on established results in the literature~\citep{peters2007reinforcement, peng2019advantage, korbak2022reinforcement, go2023aligning}, the solution that optimally maximizes reward while constraining the KL divergence, as in Eq.~\ref{eq:RL}, can be shown to have the form:

\begin{equation}\label{eq:op_policy}
    \pi_r(y\mid x) = \frac{1}{Z(x)}\pi_\text{ref}(y\mid x)\exp\left(\frac{1}{\beta}r(x, y)\right),
\end{equation}

Here, $Z(x) =\sum_{y}\pi_\text{ref}(y\mid x)\exp\left(\frac{1}{\beta}r(x, y)\right)$ acts as the normalization constant or partition function. Details of the derivation are provided in Appendix \ref{app:derivation1}. In practice, even with an MLE-based reward estimate $r_{\phi}$ approximating the ground-truth $r^*$, estimating $Z(x)$ remains computationally intensive~\citep{korbak2022reinforcement, go2023aligning}, thereby limiting the practicality of this representation. To address this, we can rearrange Eq.~\ref{eq:op_policy} and express the reward as a function of its corresponding optimal policy $\pi_r$, the baseline policy $\pi_\text{ref}$, and the still intractable $Z(\cdot)$. By taking the logarithm of both sides of Eq.~\ref{eq:op_policy} and simplifying, we arrive at:

\begin{equation}\label{eq:main_eq}
    r(x,y) =\beta \log \frac{\pi_r(y\mid x)}{\pi_\text{ref}(y\mid x)} + \beta \log Z(x).
\end{equation}

This form allows us to re-express the ground-truth reward $r^*$ in terms of its corresponding optimal policy $\pi^*$. Importantly, models such as Bradley-Terry depend only on differences in reward between output candidates, specifically ${p^*(y_1 \succ y_2 \mid x) = \sigma(r^*(x, y_1) - r^*(x, y_2))}$. Substituting the reparametrized expression from Eq.~\ref{eq:main_eq} into the preference likelihood of Eq.~\ref{eq:bradley-terry}, the partition function $Z(x)$ drops out. This yields a formulation for human preference probabilities that is solely in terms of the optimal policy $\pi^*$ and the reference policy $\pi_\text{ref}$. Thus, the optimal RLHF policy $\pi^*$ under the Bradley-Terry framework satisfies:

\begin{equation}\label{eq:objective}
    p^*(y_1\succ y_2 \mid x)=\frac{1}{1 + \exp\left(\beta \log \frac{\pi^*(y_2\mid x)}{\pi_\text{ref}(y_2\mid x)} - \beta \log \frac{\pi^*(y_1\mid x)}{\pi_\text{ref}(y_1\mid x)}\right)}
\end{equation}

A complete derivation appears in Appendix~\ref{app:derivation2}. Although this is framed using the Bradley-Terry model, analogous derivations for the broader class of Plackett-Luce models~\citep{plackett1975analysis, luce2012individual} can be found in Appendix~\ref{app:plackett_luce_models}.

With an explicit formulation for the likelihood of observed preference data in terms of the optimal policy, we can now define a maximum likelihood objective directly over a parameterized policy $\pi_\theta$. Drawing a parallel with the traditional reward modeling objective (cf. Eq.~\ref{eq:reward_model}), this leads to the following policy optimization target:

\begin{equation}\label{eq:optimum_model}
    \mathcal{L}_\text{DPO}(\pi_{\theta}; \pi_\text{ref}) = -\mathbb{E}_{(x, y_w, y_l)\sim \mathcal{D}}\left[\log \sigma \left(\beta \log \frac{\pi_{\theta}(y_w\mid x)}{\pi_\text{ref}(y_w\mid x)} - \

In this formulation, we learn an implicit reward function by adopting an alternative parameterization, resulting in an optimal policy that coincides with $\pi_\theta$. Additionally, because this method can be interpreted as training a reparameterized Bradley-Terry model, it inherits desirable theoretical properties, including consistency when the preference data is distributed under suitable conditions \cite{bong2022generalized}. We provide a more comprehensive discussion of these theoretical characteristics of DPO and its connections to related literature in Section~\ref{sec:theory}.

\textbf{How does the DPO update work?} To gain a mechanistic perspective on DPO, examining the gradient of the loss function $\mathcal{L}_\text{DPO}$ proves informative. The gradient with respect to the parameter vector $\theta$ is given by:

\begin{multline*}\label{eq:gradient}
    \nabla_\theta \mathcal{L}_\text{DPO}(\pi_\theta;\pi_\text{ref}) = \\ -\beta\mathbb{E}_{(x, y_w, y_l) \sim \mathcal{D}} \left[\underbrace{\sigma(\hat{r}_\theta(x, y_l) - \hat{r}_\theta (x, y_w))}_\text{more influence when reward assignment is inaccurate}\left[\underbrace{\nabla_\theta\log \pi(y_w \mid x)}_\text{promote $y_w$'s probability} - \underbrace{\nabla_\theta\log\pi(y_l \mid x)}_\text{reduce $y_l$'s probability}\right]\right],
\end{multline*}

where the implicit reward $\hat{r}_\theta(x, y) = \beta \log \frac{\pi_\theta(y \mid x)}{\pi_\text{ref}(y \mid x)}$ is defined in terms of the outputs of the current language model $\pi_\theta$ and a fixed reference model $\pi_\text{ref}$ (see Section~\ref{sec:theory} for details). Conceptually, the gradient update driven by $\mathcal{L}_\text{DPO}$ encourages the model to assign higher probability to the chosen (preferred) completions $y_w$, while suppressing the probabilities of the rejected (dispreferred) completions $y_l$. Crucially, the contribution of each data point is modulated by the degree to which the implicit reward function $\hat{r}_\theta$ erroneously favors the dispreferred option, modulated by the temperature parameter $\beta$. This weighting reflects the extent to which the current reward ordering conflicts with human preferences, as well as the regularization induced by the KL divergence constraint. Our experimental findings highlight the significance of this weighting scheme: omitting it, as in a simple unweighted variant, can result in catastrophic model failure (see Appendix Table~\ref{tab:unlikelihood_generations}).

\textbf{DPO overview.} The standard workflow for DPO proceeds as follows: 1) For each input prompt $x$, generate completions $y_1, y_2 \sim \pi_\text{ref}(\cdot \mid x)$, then assign human preference labels, resulting in the offline preference dataset $\mathcal{D} = \{x^{(i)}, y_w^{(i)}, y_l^{(i)}\}_{i=1}^N$; and 2) train the language model $\pi_\theta$ to minimize the objective $\mathcal{L}_\text{DPO}$ with respect to the provided $\pi_\text{ref}$, $\mathcal{D}$, and a chosen $\beta$. In practice, researchers often seek to utilize publicly available preference datasets rather than produce new completions or collect additional human labels. When such datasets have been generated with a policy $\pi^\text{SFT}$, it is preferable to initialize $\pi_\text{ref} = \pi^\text{SFT}$ if possible. In cases where $\pi^\text{SFT}$ is unavailable, $\pi_\text{ref}$ is instead initialized by fitting it to maximize the likelihood of the preferred completions $(x, y_w)$, i.e., by setting ${\pi_\text{ref} = \argmax_{\pi}\mathbb{E}_{x, y_w \sim \mathcal{D}}\left[\log \pi(y_w \mid x)\right]}$. This strategy serves to reduce the mismatch between the inaccessible true reference distribution and the operational $\pi_\text{ref}$ employed in DPO. Additional implementation and hyperparameter details are provided in Appendix~\ref{app:implementation}.

\section{Theoretical Analysis of DPO} This section delivers a deeper conceptual understanding of DPO, establishes its theoretical foundation, and discusses the method’s strengths by comparing it to actor-critic algorithms commonly applied in RLHF, such as PPO~\cite{schulman2017proximal}.

\label{sec:theory}

\subsection{Your Language Model Is Secretly a Reward Model} DPO eliminates the need to fit an explicit reward model or conduct reinforcement learning for policy optimization by employing a single maximum likelihood-based loss. Observe that the optimization problem in Eq. \ref{eq:main_eq} matches the Bradley-Terry model structure, where rewards are parameterized as $r^*(x, y) = \beta \log\frac{\pi^*_\theta(y \mid x)}{\pi_\text{ref}(y \mid x)}$. Here, optimization of the parameterized policy $\pi_{\theta}$ aligns with the reward model learning of Eq. \ref{eq:reward_model}, modulo a change of variables. The following discussion develops the theoretical reasoning behind this reparameterization, demonstrates that it does not restrict the set of possible reward functions that can be represented, and confirms the capability for exactly recovering the optimal policy. We start by introducing an equivalence relation among reward functions.

\begin{definition}
We define two reward functions $r(x, y)$ and $r'(x, y)$ to be equivalent if and only if there exists a function $f$ such that ${r(x, y) - r'(x, y) = f(x)}$.
\end{definition}

It is straightforward to verify that this definition describes an equivalence relation, partitioning the space of reward functions into equivalence classes. This leads to the following two lemmas:

\begin{lemma}\label{lemma:same_prefrence}
Within the Plackett-Luce framework—and specifically, under the Bradley-Terry model—any two reward functions belonging to the same equivalence class will generate identical preference distributions.
\end{lemma}

\begin{lemma}\label{lemma:same_policy}
    For a given constrained RL problem, all reward functions in a given equivalence class yield the same optimal policy.
\end{lemma}

We defer the relatively direct proofs to Appendix \ref{app:lemma1}. The first lemma articulates a well-documented identifiability issue associated with models in the Plackett-Luce family \cite{plackett1975analysis}. Because of this non-uniqueness, it is standard practice to incorporate additional constraints to obtain well-defined maximum likelihood estimates for the reward model in Eq. \ref{eq:reward_model} \cite{bong2022generalized}. The second lemma establishes that the optimal policy is invariant across all reward functions within an equivalence class, implying that our learning goal can be satisfied by recovering any representative from the optimal equivalence class. The following theorem, proven in Appendix~\ref{app:thm1}, makes this precise:

\begin{theorem}\label{thm:main}
    Given mild regularity conditions, all reward equivalence classes compatible with the Plackett-Luce family (and Bradley-Terry model, in particular) admit a representation of the form ${r(x, y) = \beta \log \frac{\pi(y\mid x)}{\pi_\text{ref}(y\mid x)}}$, where $\pi(y\mid x)$ denotes some policy and $\pi_\text{ref}(y \mid x)$ is a fixed reference policy.
\end{theorem}

\begin{sproof}
    Let $r(x, y)$ be any given reward function, and consider its associated optimal policy $\pi_r(y \mid x)$ defined according to Eq. \ref{eq:op_policy}. We aim to show that some reward function in the equivalence class of $r$ can be written in the form given above. Define the operator $f$ as follows:
\begin{equation}
    f(r; \pi_\text{ref}, \beta)(x, y) = r(x, y) - \beta\log\sum_{y}\pi_\text{ref}(y\mid x)\exp\left(\frac{1}{\beta}r(x, y)\right)
\end{equation}
Here, $f$ adjusts the reward by subtracting the logarithm of the partition function (with respect to $\pi_\text{ref}$), which depends solely on $x$. Thus, $f(r; \pi_\text{ref}, \beta)(x, y)$ remains in the equivalence class of $r(x, y)$. Substituting the definition of $r$ from Eq.~\ref{eq:main_eq} (valid for any reward function), it follows that $f(r; \pi_\text{ref}, \beta)(x, y) = \beta \log \frac{\pi_r(y\mid x)}{\pi_\text{ref}(y\mid x)}$. In summary, this construction yields a representative of the equivalence class in the desired form, establishing that the reparameterized reward model incurs no loss of generality.
\end{sproof}

An alternative interpretation of Theorem~\ref{thm:main} is that it characterizes precisely which reward function, among all those belonging to an equivalence class, is selected by the DPO reparameterization. Specifically, it singles out the reward function that fulfills the following condition:

\begin{equation}\label{eq:lag_p}
     \sum_{y}\underbrace{\pi_\text{ref}(y\mid x)\exp\left(\frac{1}{\beta}r(x, y)\right)}_{=\pi(y\mid x)\text{, using Thm.~\ref{thm:main} reparam.}} = 1,
\end{equation}

meaning that $\pi(y\mid x)$ forms a legitimate probability distribution with nonnegative values summing to unity. Notably, from Eq.~\ref{eq:op_policy}, we observe that Eq.~\ref{eq:lag_p} serves as the normalization (partition) function for the optimal policy arising from the reward $r(x, y)$. The essential contribution of the DPO algorithm is that by enforcing particular constraints on the otherwise indeterminate Plackett-Luce family of preference models (which includes the Bradley-Terry model), one can retain the representational power for reward models, while simultaneously ensuring that the optimal policy in Eq.~\ref{eq:op_policy} remains computationally tractable for any prompt $x$.

\subsection{Instability of Actor-Critic Algorithms} Our analytical framework also facilitates identification of failure modes in commonly used actor-critic algorithms for RLHF, such as PPO. Concentrating on the RL fine-tuning phase described in Section \ref{section:prelims}, we can relate this procedure to the control as inference perspective \cite{levine2018reinforcement} as it pertains to the constrained RL problem in \ref{eq:RL}. Assume that the policy $\pi_{\theta}(y\mid x)$ is parametrized, and the optimization seeks to minimize $\mathbb{D}_{\text{KL}}[\pi_{\theta}(y|x) \mid \mid \pi^*(y\mid x)]$, with $\pi^*$ as the optimal policy defined in Eq.~\ref{eq:optimum_model} based on the reward $r_{\phi}(y, x)$. After some algebraic manipulation, this yields the following optimization formulation:

\begin{equation}\label{eq:AC}
    \max_{\pi_{\theta}}\mathbb{E}_{\pi_{\theta}(y\mid x)}\bigg[\underbrace{r_{\phi}(x, y) -\beta\log\sum_{y}\pi_\text{ref}(y\mid x)\exp\left(\frac{1}{\beta}r_{\phi}(x, y)\right)}_{f(r_{\phi}, \pi_\text{ref}, \beta)} - \underbrace{\beta\log\frac{\pi_{\theta}(y\mid x)}{\pi_\text{ref}(y\mid x)}}_{\text{KL}}\bigg]
\end{equation}

The objective considered here aligns with that optimized in earlier research 
\citep{ziegler2020finetuning, stiennon2022learning, bai2022training, ouyang2022training}, where a DPO-equivalent reward within the reward family $r_{\phi}$ was used. In this framework, the normalization component of $f(r_{\phi}, \pi_\text{ref}, \beta)$ may be interpreted as the soft value function for the reference policy $\pi_\text{ref}$. Although the inclusion of this term does not alter the optimal solution, omitting it can result in policy gradients with elevated variance, thereby introducing instability during learning. One approach to address the normalization term is to estimate it using a learned value function; however, such optimization can be challenging. Another common strategy in preceding work involves normalizing rewards based on a human completion baseline, which serves as a one-sample Monte Carlo estimate of the normalizing term. By contrast, the DPO reparameterization produces a reward function inherently independent of any baseline.

\section{Experiments}
This section presents empirical studies assessing DPO's capacity to train policies solely from preference data. Initially, within a controlled text generation environment, we investigate how DPO balances reward maximization against KL-divergence minimization with respect to the reference policy, and compare its data efficiency to other preference learning techniques such as PPO. Subsequently, we examine DPO's effectiveness on larger-scale models and more challenging RLHF tasks, including summarization and conversational agents. Our findings indicate that, with minimal hyperparameter adjustment, DPO generally matches or surpasses leading baselines like RLHF with PPO and the method of selecting the highest-rewarded trajectory among $N$ samples from a learned reward model. The following subsections outline our experimental methodology, with further experimental details available in Appendix~\ref{app:exp_details}.

\textbf{Tasks.} We investigate three distinct open-ended text generation tasks in our experiments. Across all tasks, algorithms are trained to learn a policy using a preference dataset $\mathcal{D}=\bigl\{x^{(i)}, y_w^{(i)}, y_l^{(i)}\bigr\}_{i=1}^N$. The first task, \textbf{controlled sentiment generation}, involves $x$ as an initial segment of a movie review taken from the IMDb dataset \cite{maas-EtAl:2011:ACL-HLT2011}, with the requirement that the model generate $y$ that conveys positive sentiment. To enable systematic evaluation, for this task we synthetically generate pairs of outputs using a pre-trained sentiment classifier, selecting pairs such that $p(\text{positive}\mid x,y_w)>p(\text{positive}\mid x,y_l)$. For supervised fine-tuning (SFT), we adapt GPT-2-large through continued training on positive reviews from the IMDB training set, following details elaborated in App~\ref{app:sentiment_details}. 

In the second task, \textbf{summarization}, $x$ represents a forum post sourced from Reddit, and the model's objective is to produce a summary $y$ highlighting the main ideas. We use the Reddit TL;DR dataset for summarization \citep{volske-etal-2017-tl}, supplementing it with human preference data from \citeauthor{stiennon2022learning}, in line with previous studies. For this task, the SFT model is trained on summaries of Reddit forum posts written by humans,\footnote{\url{https://huggingface.co/CarperAI/openai_summarize_tldr_sft}} utilizing the TRLX \citep{leandro_von_werra_2023_7790115} library for RLHF training. The preference annotations, which were created by \citeauthor{stiennon2022learning}, come from outputs of a different yet comparably fine-tuned SFT model.

The third task, \textbf{single-turn dialogue}, presents $x$ as a user prompt, varying from scientific inquiries to advice-seeking, for which the system must return an informative and engaging reply $y$. For this purpose, we employ the Anthropic Helpful and Harmless dialogue dataset \citep{bai2022training}, comprising 170k human-assistant conversations. Each dialog concludes with two responses produced by a (large, unspecified) language model, along with a human annotation indicating the preferred reply. As there is no pre-existing SFT model in this scenario, we construct one by fine-tuning a generic language model using only the preferred responses.

\textbf{Evaluation.} We adopt two complementary evaluation strategies in our experimental framework. For the controlled sentiment generation context, we quantitatively assess each algorithm’s performance with respect to the constrained reward maximization objective by plotting its trade-off frontier between achieved reward and KL-divergence from the reference policy; this is feasible here due to the availability of the true reward function (provided by a sentiment classifier). In contrast, such explicit reward supervision is unavailable in practical, real-world applications. Consequently, we turn to measuring each algorithm’s \textit{win rate} versus a baseline policy, where GPT-4 serves as an automated stand-in for human evaluators to rate summary quality in the summarization task and helpfulness in the single-turn dialogue task. For summarization, test set reference summaries act as baselines; for dialogue, the baseline consists of the dataset’s preferred human response. While previous literature reports that LMs outperform conventional metrics as automatic judges \citep{Chen2023ExploringTU}, we bolster our choice of GPT-4 for evaluation by conducting a human study (see Sec.~\ref{sec:human-judgments}), where we observe a high correlation between GPT-4’s assessments and those of human annotators, often rivaling or surpassing inter-annotator agreement levels.

\textbf{Methods.} Beyond DPO, our empirical comparison includes multiple established approaches for aligning language models with human preferences. For the summarization task, we use zero-shot prompting of \textbf{GPT-J} \citep{gpt-j}, while in the dialogue setting, we employ 2-shot prompting using \textbf{Pythia-2.8B} \citep{biderman2023pythia}. We further assess an \textbf{SFT} model alongside a \textbf{Preferred-FT} variant, the latter trained via supervised fine-tuning on chosen completions $y_w$—sourced either from the SFT model (in the controlled sentiment and summarization domains) or from a generic LM (for single-turn dialogue). Among alternative pseudo-supervised techniques, we evaluate \textbf{Unlikelihood}~\citep{welleck2019neural}, where optimization targets increasing the likelihood of $y_w$ while penalizing the likelihood assigned to $y_l$, moderated by an optional scaling factor $\alpha\in[0,1]$ applied to the unlikelihood loss. Additionally, we include \textbf{PPO} \citep{schulman2017proximal}, which utilizes a reward model trained on preference data, and \textbf{PPO-GT}, an oracle leveraging access to the true reward function in the controlled sentiment setup. For sentiment experiments, we benchmark two versions of PPO-GT: a standard off-the-shelf implementation \cite{leandro_von_werra_2023_7790115} and a custom variant with normalized rewards and refined hyperparameter tuning—improvements also applied to PPO when using learned rewards. Finally, we introduce a \textbf{Best of $N$} strategy, which involves sampling $N$ completions from either the SFT model or Preferred-FT (in dialogue), and selecting the response with the highest score as per the learned reward model. Although highly effective, this method dissociates reward model fidelity from PPO-based optimization, and remains computationally expensive for any substantial $N$, given that it demands $N$ samples per test query.

\subsection{How well can DPO optimize the RLHF objective?}

In standard RLHF methodologies, the KL-constrained reward maximization objective serves to simultaneously maximize reward acquisition and prevent substantial divergence from the reference policy. Consequently, when assessing different approaches, it is important to consider both the attained reward and the magnitude of KL divergence; a marginal increase in reward accompanied by a substantial rise in KL is often not preferable. Figure~\ref{fig:frontier-tldr-main} illustrates the reward-KL tradeoff achieved by a range of algorithms under the sentiment analysis configuration. For each method, we conduct several independent training sessions, varying a conservativeness-related hyperparameter across runs (target KL values $\in\{3,6,9,12\}$ for PPO, $\beta \in \{0.05,0.1,1,5\}$, $\alpha\in\{0.05,0.1,0.5,1\}$ for unlikelihood, and altering random seeds for preferred-FT). This comprehensive sweep totals 22 experiments. Every 100 steps during training, up to convergence, we assess the learned policies on a suite of test prompts, recording both the mean reward (according to the ground-truth reward model) and the mean sequence-level KL\footnote{That is, the sum of the per-timestep KL-divergences.} relative to the reference policy, $\text{KL}\left(\pi\mid \mid \pi_\text{ref}\right)$. Our analysis reveals that DPO yields the most effective reward-KL frontier by a significant margin, securing maximal rewards while maintaining minimal KL divergence. This finding stands out for several reasons. Notably, although DPO and PPO both optimize for the same underlying objective, DPO is far more efficient; the DPO tradeoff curve between reward and KL is strictly better than that of PPO. Moreover, DPO surpasses PPO even in situations where PPO is provided access to exact reward labels (PPO-GT).

\subsection{Can DPO scale to real preference datasets?}
\label{sec:dpo-real-datasets}
We proceed to assess the effectiveness of DPO fine-tuning using tasks in summarization and single-turn dialogue. For summarization, it is well-known that automatic evaluation measures such as ROUGE frequently exhibit low alignment with human judgment~\citep{stiennon2022learning}, and previous studies have shown that fine-tuning language models with PPO based on human feedback leads to superior summarization performance. Our comparative evaluation consists of generating completions from each method on the test portion of the TL;DR summarization dataset, followed by calculating the mean win rate against the set of reference completions. All samples are generated at a range of temperature values between 0.0 and 1.0, and the corresponding win rates are illustrated in Figure~\ref{fig:frontier-tldr-main} (right panel). The DPO, PPO, and Preferred-FT approaches all further train the identical GPT-J SFT model\footnote{\url{https://huggingface.co/CarperAI/openai_summarize_tldr_sft}}. Our results demonstrate that DPO reaches an approximate win rate of 61\% at a sampling temperature of 0.0, surpassing PPO, which achieves about 57\% at its optimal sampling temperature of 0.0. Additionally, DPO attains a greater maximum win rate than the best of $N$ baseline. It should be noted that minimal hyperparameter tuning was applied to DPO's $\beta$ value, indicating that these outcomes may not represent the full capability of DPO. Furthermore, DPO maintains consistently higher robustness to temperature fluctuations than PPO, which experiences considerable performance loss at elevated temperatures, regressing towards the baseline GPT-J model. Preferred-FT, in contrast, exhibits negligible gains over the original SFT model. Direct human assessments comparing DPO and PPO (see Section~\ref{sec:human-judgments}) further confirm these findings: at a temperature of 0.25, DPO-generated samples were favored over those from PPO 58\% of the time.

For single-turn dialogues, we assess the various approaches on the subset of the Anthropic HH dataset's test split \citep{bai2022training} that involves one round of interaction between human and assistant. In GPT-4-based evaluations, we determine win rates for each method by referencing the test set's preferred completions. Given the absence of a canonical SFT model for this benchmark, we initialize with the pre-trained Pythia-2.8B, apply Preferred-FT to fine-tune a reference model on the preferred completions to ensure outputs remain in-distribution, and subsequently train with DPO. Additionally, we benchmark against the strongest performer among 128 Preferred-FT completions (as performance plateaus beyond 128 completions on this dataset; refer to Appendix Figure~\ref{fig:best-of-n}) and a 2-shot prompted variant of the Pythia-2.8B base model, observing that DPO achieves equal or superior results when each method is tuned to its optimal temperature. Furthermore, an RLHF model—trained via PPO on the Anthropic HH dataset \footnote{\url{https://huggingface.co/reciprocate/ppo_hh_pythia-6B}} and sourced from a widely recognized implementation \footnote{\url{https://github.com/CarperAI/trlx/tree/main/examples/hh}}—is evaluated, but no combination of prompts or temperatures yields improvements over the original Pythia-2.8B base. Considering both TL;DR outcomes and the fact that PPO and Best of 128 optimize the same reward signal, we treat the latter as an approximate indicator of PPO's performance ceiling. In summary, DPO stands out as the only scalable technique that surpasses the Anthropic HH dataset's preferred completions, matching or exceeding the effectiveness of the computationally intensive Best of 128. Figure~\ref{fig:dialogue-main} further illustrates that DPO attains peak performance in comparatively few iterations.

\subsection{Generalization to a new input distribution}

To further investigate the robustness of PPO and DPO under distribution shifts, we assess both policies—previously trained on the Reddit TL;DR summarization task—on a distinct out-of-distribution dataset: the test split of CNN/DailyMail, which consists of news articles \citep{nallapati-etal-2016-abstractive}. We utilize the optimal sampling temperatures determined from the TL;DR experiment (0 and 0.25). The performance metrics are shown in Table~\ref{tab:ood}. GPT-4’s win rate over the reference summaries was measured using an evaluation prompt identical to the one for Reddit TL;DR, except the phrase ``forum post'' was switched to ``news article''. Across this novel input distribution, the DPO policy consistently surpasses PPO by a considerable margin. These findings suggest that DPO, despite lacking access to additional unlabeled Reddit TL;DR prompts utilized by PPO, exhibits generalization abilities that match or exceed those of PPO.

\subsection{Validating GPT-4 judgments with human judgments}
\label{sec:human-judgments}

We perform a human evaluation to examine the correspondence between GPT-4’s judgments and those of human raters, drawing upon the outputs from the TL;DR summarization experiment and employing two distinct GPT-4 prompting strategies. The \textbf{GPT-4 (S)} (simple) prompt instructs GPT-4 to select the summary that more effectively conveys the post’s important points. The \textbf{GPT-4 (C)} (concise) prompt further directs the model to prefer more succinct summaries; this prompt is included because with the simple prompt, GPT-4 displays a bias towards longer and more repetitive summaries compared to human preferences. Full details of both prompts are available in Appendix~\ref{app:prompts}. Our study examines three system comparisons—one with the top-performing (DPO, temp. 0.25), one with the lowest-performing (PPO, temp. 1.0), and an intermediate-quality system (SFT, temp. 0.25)—all evaluated against PPO using greedy sampling (its best setting). This sampling strategy is designed to capture a broad spectrum of sample qualities. Our results indicate that for both prompting approaches, GPT-4's judgments align with those of human evaluators at rates comparable to inter-human agreement, implying GPT-4 can serve as a reliable stand-in for human assessment (note that due to a limited pool of raters, multiple human annotations were only gathered for DPO and PPO-1 comparisons). Overall, the \textbf{GPT-4 (C)} prompt tends to yield win rates that better reflect human evaluations; as a result, we base our primary analyses in Section~\ref{sec:dpo-real-datasets} on this prompt. For further information regarding the human evaluation process, the rater interface, and the participant list, see Appendix~\ref{app:human-study}.

\section{Discussion}
Preference-based learning represents an effective and extensible approach for developing proficient and value-aligned language models. In this work, we have presented DPO, a straightforward methodology for preference-driven language model training that bypasses the need for reinforcement learning techniques. Instead of adapting the preference learning task to fit conventional RL paradigms simply to apply standard RL algorithms, DPO constructs a correspondence between language model policies and reward functions, allowing language models to be trained to adhere to human preferences \textit{directly} using only a cross-entropy loss, without resorting to reinforcement learning or sacrificing generality. Remarkably, DPO achieves comparable or superior results to prevailing RLHF methods, such as those using PPO, all while requiring minimal hyperparameter adjustment, thereby significantly lowering the obstacles to training language models from human preferences.

\textbf{Limitations \& Future Work.} The findings presented prompt a series of avenues for further investigation. It remains to be seen how policies optimized via DPO extend to out-of-distribution data relative to those trained with explicit reward functions. Preliminary evidence indicates that DPO-trained policies generalize on par with those refined by PPO, though more thorough analyses are warranted. Further, questions remain regarding whether the use of self-generated labels from a DPO-trained policy can enable effective leveraging of prompts lacking human annotation. Additionally, the nature and extent of reward over-optimization in direct preference optimization merit careful study; for instance, the minor performance dip observed in Figure~\ref{fig:dialogue-main}-right may exemplify this issue. Moreover, while our evaluations are confined to models with up to 6B parameters, scaling DPO to significantly larger, state-of-the-art architectures is a promising trajectory for future research. On the evaluation front, our results show that win rates measured via GPT-4 are prompt-sensitive, suggesting that refining automated evaluation protocols could be a productive line of inquiry. Ultimately, the scope of DPO’s applicability extends beyond preference-based language model training and includes a wide array of generative modeling problems across various modalities.

\end{document}